% ============================================================
\chapter{M\'ethodes It\'eratives pour Syst\`emes Lin\'eaires}
\label{ch:iteratives}
% ============================================================

\section{Exemple motivant}
En \'el\'ements finis, la matrice $A$ est creuse ($\BigO(n)$ coefficients non nuls sur $n^2$). LU co\^ute $\BigO(n^3)$, inacceptable pour $n=10^6$. Les m\'ethodes it\'eratives exploitent la structure creuse : chaque it\'eration co\^ute $\BigO(n)$.

\section{Principe g\'en\'eral}

\begin{definition}[M\'ethode de d\'ecomposition]
On \'ecrit $A=M-N$ avec $M$ inversible facile. L'it\'eration est :
\[
Mx^{(k+1)} = Nx^{(k)}+b, \quad\text{soit}\quad x^{(k+1)}=Gx^{(k)}+c,
\]
o\`u $G=M^{-1}N$ est la \textbf{matrice d'it\'eration} et $c=M^{-1}b$.
\end{definition}

\begin{theoreme}[Convergence]\label{thm:iter-conv}
La m\'ethode converge pour tout $x^{(0)}$ si et seulement si $\rho(G)<1$ (rayon spectral).
\end{theoreme}

\begin{proof}
$e^{(k)}=x^{(k)}-x^*=G^k e^{(0)}$. Or $\|G^k\|\to 0 \iff \rho(G)<1$.
\end{proof}

\section{M\'ethode de Jacobi}\index{Jacobi}

\begin{definition}
$A=D-E-F$ (diagonale, triangulaire inf., triangulaire sup.). $M=D$ :
\[
x_i^{(k+1)} = \frac{1}{a_{ii}}\left(b_i-\sum_{j\neq i}a_{ij}x_j^{(k)}\right).
\]
Matrice d'it\'eration : $G_J=D^{-1}(E+F)$.
\end{definition}

\section{M\'ethode de Gauss--Seidel}\index{Gauss--Seidel}

\begin{definition}
$M=D-E$ :
\[
x_i^{(k+1)} = \frac{1}{a_{ii}}\left(b_i-\sum_{j<i}a_{ij}x_j^{(k+1)}-\sum_{j>i}a_{ij}x_j^{(k)}\right).
\]
Matrice d'it\'eration : $G_{GS}=(D-E)^{-1}F$.
\end{definition}

\begin{theoreme}[Convergence pour matrices SDP]
Si $A$ est SDP, Gauss--Seidel converge.
\end{theoreme}

\begin{theoreme}[Convergence pour matrices \`a diagonale dominante]
Si $A$ est \`a diagonale strictement dominante ($|a_{ii}|>\sum_{j\neq i}|a_{ij}|$), alors Jacobi et Gauss--Seidel convergent.
\end{theoreme}

\section{Relaxation (SOR)}\index{SOR}

\begin{definition}
On introduit un param\`etre $\omega\in(0,2)$ :
\[
x_i^{(k+1)} = (1-\omega)x_i^{(k)} + \frac{\omega}{a_{ii}}\left(b_i-\sum_{j<i}a_{ij}x_j^{(k+1)}-\sum_{j>i}a_{ij}x_j^{(k)}\right).
\]
$\omega=1$ : Gauss--Seidel. $\omega>1$ : sur-relaxation. $\omega<1$ : sous-relaxation.
\end{definition}

\begin{theoreme}
SOR converge ssi $0<\omega<2$. L'$\omega$ optimal d\'epend de $\rho(G_J)$.
\end{theoreme}

\section{M\'ethode du gradient conjugu\'e}\index{gradient conjugue@gradient conjugu\'e}

\begin{definition}
Pour $A$ SDP, r\'esoudre $Ax=b$ \'equivaut \`a minimiser $\phi(x)=\frac{1}{2}x^\top Ax-b^\top x$.
\end{definition}

\begin{algorithm}[H]
\caption{Gradient conjugu\'e}\label{alg:cg}
\KwEntree{$A$ SDP, $b$, $x_0$, tol\'erance $\eps$}
\KwSortie{$x\approx A^{-1}b$}
$r_0\gets b-Ax_0$, $p_0\gets r_0$\;
\Pour{$k=0,1,2,\ldots$}{
  $\alpha_k\gets\frac{r_k^\top r_k}{p_k^\top Ap_k}$\;
  $x_{k+1}\gets x_k+\alpha_k p_k$\;
  $r_{k+1}\gets r_k-\alpha_k Ap_k$\;
  \Si{$\|r_{k+1}\|<\eps$}{\Return{$x_{k+1}$}}
  $\beta_k\gets\frac{r_{k+1}^\top r_{k+1}}{r_k^\top r_k}$\;
  $p_{k+1}\gets r_{k+1}+\beta_k p_k$\;
}
\end{algorithm}

\begin{theoreme}[Convergence du gradient conjugu\'e]
En arithm\'etique exacte, CG converge en au plus $n$ it\'erations. L'erreur v\'erifie :
\[
\|e_k\|_A \leq 2\left(\frac{\sqrt{\kappa}-1}{\sqrt{\kappa}+1}\right)^k\|e_0\|_A, \quad \kappa=\cond_2(A).
\]
\end{theoreme}

\section{Impl\'ementation}

\begin{codeblock}[title=\textbf{Python}]
\begin{minted}{python}
import numpy as np

def jacobi(A, b, x0, tol=1e-10, maxiter=1000):
    D = np.diag(np.diag(A))
    R = A - D
    x = x0.copy()
    for k in range(maxiter):
        x_new = np.linalg.solve(D, b - R @ x)
        if np.linalg.norm(x_new - x) < tol:
            return x_new, k+1
        x = x_new
    return x, maxiter

def conjugate_gradient(A, b, x0, tol=1e-10, maxiter=1000):
    x = x0.copy()
    r = b - A @ x
    p = r.copy()
    rs_old = r @ r
    for k in range(maxiter):
        Ap = A @ p
        alpha = rs_old / (p @ Ap)
        x += alpha * p
        r -= alpha * Ap
        rs_new = r @ r
        if np.sqrt(rs_new) < tol:
            return x, k+1
        p = r + (rs_new / rs_old) * p
        rs_old = rs_new
    return x, maxiter

# Test : matrice tridiagonale
n = 100
A = np.diag(2*np.ones(n)) - np.diag(np.ones(n-1), 1) - np.diag(np.ones(n-1), -1)
b = np.ones(n)
x0 = np.zeros(n)

x_j, nj = jacobi(A, b, x0)
x_cg, ncg = conjugate_gradient(A, b, x0)
print(f"Jacobi: {nj} iterations")
print(f"CG:     {ncg} iterations")
\end{minted}
\end{codeblock}

\begin{codeblock}[title=\textbf{Julia}]
\begin{minted}{julia}
function cg(A, b, x0; tol=1e-10, maxiter=1000)
    x = copy(x0)
    r = b - A * x
    p = copy(r)
    rs = dot(r, r)
    for k in 1:maxiter
        Ap = A * p
        alpha = rs / dot(p, Ap)
        x .+= alpha .* p
        r .-= alpha .* Ap
        rs_new = dot(r, r)
        sqrt(rs_new) < tol && return x, k
        p .= r .+ (rs_new / rs) .* p
        rs = rs_new
    end
    return x, maxiter
end
\end{minted}
\end{codeblock}

\section{Exercices}

\begin{exercice}[$\star$]
Appliquer 3 it\'erations de Jacobi et Gauss--Seidel au syst\`eme $\begin{pmatrix}4&1\\1&3\end{pmatrix}x=\begin{pmatrix}1\\2\end{pmatrix}$, $x^{(0)}=(0,0)^\top$.
\end{exercice}

\begin{exercice}[$\star\star$]
Montrer que pour une matrice tridiagonale $T_n$ de taille $n$, $\rho(G_J)=\cos(\pi/(n+1))$. En d\'eduire le nombre d'it\'erations de Jacobi pour une pr\'ecision $\eps$.
\end{exercice}

\begin{exercice}[$\star\star\star$ -- Projet]
Comparer Jacobi, Gauss--Seidel, SOR et CG pour la discr\'etisation 2D du Laplacien $-\Delta u=f$ sur $[0,1]^2$. Tracer le nombre d'it\'erations en fonction de $n$.
\end{exercice}

\begin{formules}
\begin{align*}
\text{Jacobi} &: x^{(k+1)}=D^{-1}(b-(A-D)x^{(k)}) \\
\text{Gauss--Seidel} &: G=(D-E)^{-1}F \\
\text{CG} &: \|e_k\|_A\leq 2\left(\frac{\sqrt{\kappa}-1}{\sqrt{\kappa}+1}\right)^k\|e_0\|_A
\end{align*}
\end{formules}
